\question[3]
Verschlüsselung begegnet uns im Alltag an vielen Stellen, auch wenn 
sie oft im Hintergrund abläuft und nicht bewusst wahrgenommen wird.

Nenne \textbf{drei} unterschiedliche Anwendungsbereiche, in denen 
Verschlüsselung eingesetzt wird, und beschreibe jeweils kurz, welche 
Daten dabei geschützt werden und wovor.

\begin{solutionbox}{9cm}

\begin{itemize}
    \item \textbf{Verschlüsselte Kommunikation (z.\,B.\ Messenger wie 
    Signal oder WhatsApp):} Die Inhalte von Nachrichten werden 
    verschlüsselt, sodass nur Sender und Empfänger sie lesen können. 
    Geschützt wird der Inhalt der Kommunikation vor Mitlesen durch 
    Dritte, etwa den Anbieter des Dienstes oder Angreifer im Netz.
    
    \item \textbf{Datenübertragung im Web (HTTPS):} Die Verbindung 
    zwischen Browser und Webserver wird verschlüsselt (z.\,B.\ beim 
    Online-Banking oder Onlineshopping). Geschützt werden dabei 
    übertragene Daten wie Passwörter, Zahlungsinformationen oder 
    persönliche Angaben vor dem Mitlesen oder Verändern durch Dritte 
    während der Übertragung.
    
    \item \textbf{Verschlüsselte Speicherung von Daten 
    (z.\,B.\ Festplattenverschlüsselung, Cloud-Speicher):} Dateien 
    werden auf dem Gerät oder in der Cloud nur in verschlüsselter Form 
    gespeichert. Geschützt werden die gespeicherten Daten vor unbefugtem 
    Zugriff, etwa bei Diebstahl des Geräts oder Einbruch in den 
    Cloud-Speicher.
    
    \item \textbf{WLAN-Verschlüsselung (z.\,B.\ WPA2/WPA3):} Der 
    Datenverkehr zwischen Gerät und WLAN-Router wird verschlüsselt. 
    Geschützt werden übertragene Daten vor dem Mitlesen durch andere 
    Personen in Reichweite des Funknetzes.
\end{itemize}

\end{solutionbox}


